% Part: lambda-calculus
% Chapter: introduction
% Section: lambda-computable

\documentclass[../../../include/open-logic-section]{subfiles}

\begin{document}

\olfileid[pt]{lam}{int}{cmp}
\olsection{Funções \usetoken{S}{lambda definable} são computáveis}

\begin{thm}
\ollabel{thm:lambda-computable}
Se uma função parcial $f$ é !!{lambda defined} por um termo lambda, então é
computável.
\end{thm}

\begin{proof}
Suponha que uma função~$f$ seja !!{lambda defined} por um termo lambda~$X$.
Descreveremos um procedimento informal para calcular~$f$. Com entrada
$m_0$, \dots,~$m_{n-1}$, escreva o termo $X \num m_0 \ldots \num
m_{n-1}$. Construa uma árvore: primeiro, escreva todas as reduções em um
passo do termo original; abaixo delas, escreva todas as reduções em um passo
desses termos (isto é, as reduções em dois passos do termo original); e
continue. Se em algum momento você chegar a um numeral, devolva-o como
resposta; caso contrário, a função não está definida.

Um apelo à tese de Church nos diz que essa função é computável. Uma maneira
melhor de provar o teorema seria dar uma descrição recursiva desse
procedimento de busca. Por exemplo, poderíamos definir uma sequência de
funções e relações recursivas primitivas: ``$\fn{IsASubterm}$'',
``$\fn{Substitute}$'', ``$\fn{ReducesToInOneStep}$'',
``$\fn{ReductionSequence}$'', ``$\fn{Numeral}$'' etc. O procedimento
recursivo parcial para calcular $f(m_0, \dots, m_{n-1})$ consiste então em
procurar uma sequência de reduções em um passo que comece em
$X \num{m_0} \dots \num{m_{n-1}}$ e termine em um numeral, e devolver o
número correspondente a esse numeral. Os detalhes são longos e tediosos,
mas, fora isso, rotineiros.
\end{proof}

\end{document}

